%% P10 --- Wartości własne, formy kwadratowe i dodatnia określoność
%%
%% Zalążek, wygenerowany przez tools/gen_stubs.py z tools/programs.json.
%% Poniższy opis jest kontraktem tego programu: co ma uzasadnić i co ma dać
%% czytelnikowi. Napisanie programu oznacza usunięcie bloku \programstub{} i
%% zastąpienie go programem -- po czym ten plik nie jest już generowany.

\program{Wartości własne, formy kwadratowe i dodatnia określoność}
\label{prog:P10}

\programstub{%
\textit{[Opis do przetłumaczenia --- patrz wydanie angielskie.]}

Argument: an eigenvector is a direction the matrix only stretches; a
symmetric matrix has a full orthogonal set of them (the spectral theorem,
stated and used, not proved); a quadratic form is a bowl, a saddle or a
ridge, and its eigenvalues say which. Payoff: a covariance matrix is
symmetric and positive semi-definite, so PCA is an eigen-decomposition and
not a black box; the spectral norm as the largest stretch, which is the
Lipschitz constant that bounds how much a layer can amplify; **the Hessian's
eigenvalues are the shape of the loss basin**, which is where
\enquote{ravine}, \enquote{sharp minimum} and \enquote{the learning rate is
too high} all become one statement \dash{} collected here and spent in P17
and P20.

FORWARD REFERENCE, declared: this program's payoff uses the covariance
matrix, defined in P24. State the two facts it needs -- that a covariance
matrix is symmetric and positive semi-definite -- in the Learning outcomes
with the pointer, exactly as P21 does for its probability prerequisites.

\medskip
\textbf{Planowana długość:} 65 ramek.
}
